\section{问题三的预报更新与购电调整}
\subsection{调整交易的结算与经济含义}
问题三允许在6:00、12:00和18:00获得新的光伏预报，并据此调整尚未交付的购电量。其作用不仅是提高储能调度对供需变化的适应能力，还包括提前修正购电不足或过量的安排。评价这些调整时，需要同时计入紧急补电、取消计划和新增购电的费用。

记$g_t$为0:00制定的原计划，$a_t$为第$t$段实际交付的最终有效购电量。本文假设取消部分只支付原价的50\%作为违约费，即退回取消部分的原购电费后收取违约金；新增部分按1.5倍电价购买。定义
\begin{equation}
\delta_t^+=(a_t-g_t)^+,\qquad \delta_t^-=(g_t-a_t)^+,
\end{equation}
则当天实际费用为
\begin{align}
C_d&=\sum_t p_tg_t-\sum_t p_t\delta_t^-+\sum_t0.5p_t\delta_t^-
       +\sum_t1.5p_t\delta_t^++\sum_t5p_te_t\\
&=\sum_t\left[p_tg_t-0.5p_t\delta_t^-+1.5p_t\delta_t^++5p_te_t\right].
\end{align}
例如，电价为1元/kWh，某段原计划100 kWh，最终减至80 kWh，其非紧急费用为90元。若先改为80 kWh、后改为90 kWh，则按最终数量结算95元。每段始终相对0点原计划计算一次偏差，不对历次预报版本重复计费。

取消的电量不再交付，不能继续供给负载或给电池充电；已经交付但未利用的电量仍收费。实际能量平衡采用
\begin{equation}
a_t+P_t+d_t+e_t=L_t+c_t+w_t.
\end{equation}
储电量递推、容量边界、功率上限及单程效率均与问题二相同。通过非负变量$\delta_t^+,\delta_t^-$及等式$a_t-g_t=\delta_t^+-\delta_t^-$可将结算函数线性化。同时增加二者会使费用上升，因此费用最小解不会同时出现无意义的增购和减购。

\subsection{滚动预报与十分钟电量映射}
附件三在每日四个发布时刻提供随后24个整点的光伏功率预报。设发布时间为$\tau$，第$j$个点对应$\tau+j$小时。以刚结束区间的已知光伏功率作为$\tau$处的起点，对相邻整点进行分段线性插值。若十分钟区间两端的插值功率为$\widehat P_t^{\mathrm{left}}$和$\widehat P_t^{\mathrm{right}}$，则该段预测电量为
\begin{equation}
\widehat V_t=\frac{\widehat P_t^{\mathrm{left}}+\widehat P_t^{\mathrm{right}}}{2}\times\frac{1}{6}.
\end{equation}
该转换保留整点功率信息，并对所假定的线性变化进行积分。18:00发布的第6个点对应次日0:00，跨日预报也参与储能前瞻。

规划终点仍为次日24:00。对当前24小时预报覆盖不到的远端时段，使用最近完整一天对应时段的光伏电量作补充预测；负载沿用上周同一时段预测。补充部分仅用于前瞻，不视为新发布的预报。评价期和初始状态仍为2月1日至12月31日、6000 kWh，实际储电量连续跨日传递。

对每个发布时间分别重建历史预测误差：在历史起点只使用其当时可见的实测与对应预报，按上述规则补足规划范围，再与随后两日实测比较。采用最近28个完整历史样本及7个代表情景，保留配对的负载、光伏轨迹和净负荷极端。采用新预报后重新构造情景；本次发布以来的实测误差用于更新概率，保留5\%原始权重。这样既避免使用尚未发布的预报，也不将上一版本的误差直接当作新版本误差。

\subsection{日前计划、日内调整与实时执行}
0:00依据当时预报建立问题二形式的多情景日前模型，确定当天144段原计划。日内采用新预报时，给定实际储电量和已执行区间，对当天剩余区间的$a_t$共同优化，各情景允许不同的储能与紧急补电。目标包括当天剩余交易的结算费用、期望紧急购电费，以及次日前瞻购电费；远端储电目标为6000 kWh。

每次优化都保留当前有效购电计划作为可行候选。在相同情景、状态和末端条件下，分别计算保持计划与允许调整的最优预计费用；若改善不超过0.01元，则保持当前计划。否则在最优费用数值容差内，优先选取总变更电量较小的方案。该判断使用已经包含调整费用的目标值，不再额外扣除一次违约费或增购费。

在两个预报发布时间之间，每十分钟固定当前有效计划，更新观测与实际储电量，重新优化储能并只执行当前共同动作。当前供给有盈余时只允许充电或舍弃；存在缺额时只允许放电或紧急补电。次日购电及未来情景动作仅作为前瞻，随后重新决策。由于没有在0点预先构造未来所有预报版本的决策树，本方法是分时修正策略，而不是完整多阶段随机规划的全局最优解。

\subsection{全年对照与预报时刻的价值}
对6、12、18点预报的全部八种选用组合进行独立连续回测，0点预报始终采用。另设接受全部预报但仅调整储能的策略，以区分信息更新和修改购电量的作用。各策略均保留不调整的选择，所有费用由最终交付量与实际紧急购电结算。

\begin{table}[htbp]\centering\scriptsize
\caption{问题三各策略的334天实际运行结果}\label{tab:q3compare}
\resizebox{\textwidth}{!}{\begin{tabular}{lrrrrr}
\toprule
策略 & 费用/万元 & 紧急电量/kWh & 天数 & 未利用电量/kWh & 期末储电/kWh \\
\midrule
仅0点预报 & 1468.12 & 301147.22 & 208 & 2339713.05 & 8314.11 \\
6点调整 & 1424.68 & 196735.01 & 186 & 1983450.36 & 8314.11 \\
12点调整 & 1407.42 & 192857.14 & 195 & 1871994.03 & 8359.44 \\
18点调整 & 1422.88 & 197939.78 & 208 & 2327996.23 & 8486.34 \\
6、12点调整 & 1396.07 & 166577.57 & 186 & 1738038.20 & 8359.44 \\
6、18点调整 & 1401.84 & 145568.06 & 188 & 1978655.83 & 8493.72 \\
12、18点调整 & 1394.77 & 163326.91 & 195 & 1868391.18 & 8493.72 \\
6、12、18点调整 & 1382.99 & 136287.62 & 186 & 1734278.82 & 8493.72 \\
全预报仅调储能 & 1468.55 & 307111.99 & 218 & 2353632.40 & 8314.11 \\
原计划照付口径 & 1400.86 & 127399.48 & 166 & 2439450.28 & 8493.72 \\
\bottomrule
\end{tabular}}
\end{table}
三个时刻均采用的主方案总费用为\QthreeCost\ 万元，紧急购电量为\QthreeEmergency\ kWh，发生于\QthreeDays\ 天。以仅采用0点预报为基准，费用减少\QthreeDifference\ 元，降幅为\QthreeSaving\%。费用构成为原计划金额\QthreePlanCost\ 万元，扣除取消部分退款\QthreeRefund\ 万元，加上违约费\QthreePenalty\ 万元、增购费\QthreeExtra\ 万元和紧急购电费\QthreeEmergencyCost\ 万元。

仅更新预报并调整储能，相对0点基准的实际费用增加4284.87元。进一步允许修改购电量，相对仅更新储能的方案，费用减少855533.38元。本次对照中，购电修订提供的灵活性是费用改善的主要来源。这是不同连续运行策略之间的费用差，包含其后续储电状态变化的影响。
在八种预报组合中，本次回测费用最低的是“6、12、18点调整”，总费用为1382.9943万元。该排序是样本内的运行比较，不作为重新选择主方案参数的依据。在其余预报组合固定时，加入6点预报的全年费用减少量介于113525.26至434431.22元。在其余预报组合固定时，加入12点预报的全年费用减少量介于188471.60至606987.13元。在其余预报组合固定时，加入18点预报的全年费用减少量介于126510.18至452429.31元。因此，在本次数据与模型设定下，采用6、12、18点预报并允许购电调整具有费用优势；接收预报不意味着每次都必须修改计划。在原计划照付的另一结算口径下，全部预报方案相对无调整基准的费用减少672587.38元，全年减购量为0.0000 kWh。

\begin{figure}[p]\centering
\includegraphics[width=\textwidth]{figures/q3_compare.pdf}
\caption{预报选用及调整策略的费用与紧急购电比较。}
\end{figure}
\begin{figure}[p]\centering
\includegraphics[width=\textwidth]{figures/q3_year.pdf}
\caption{主方案逐月购电费用构成。净基础费用为原计划金额减去取消部分退款。}
\end{figure}

\subsection{指定日期与敏感性分析}
表\ref{tab:q3selected}至表\ref{tab:q3emergency}给出四个指定日期的计划、最终调整量及实际储能运行。调整后电量指非紧急交付量，加上紧急电量才是总购电量。
\begin{table}[htbp]\centering\scriptsize
\caption{指定日期的电量、总费用及首尾储电量（电量单位：kWh）}\label{tab:q3selected}
\resizebox{\textwidth}{!}{\begin{tabular}{lrrrrrr}
\toprule
日期 & 原计划电量 & 调整后电量 & 紧急电量 & 总费用/元 & 日初储电 & 日末储电 \\
\midrule
03-20 & 65824.81 & 61938.11 & 476.95 & 42591.41 & 8467.52 & 4646.88 \\
06-21 & 42586.35 & 42325.44 & 0.00 & 24691.61 & 5592.84 & 10670.88 \\
09-23 & 74771.47 & 68804.75 & 455.94 & 47093.43 & 7716.50 & 8166.23 \\
12-21 & 98420.21 & 96302.01 & 0.00 & 63767.32 & 8177.36 & 8256.55 \\
\bottomrule
\end{tabular}}
\end{table}
\begin{table}[htbp]\centering\scriptsize
\caption{指定时段的原计划与最终调整电量（kWh）}\label{tab:q3purchase}
\resizebox{\textwidth}{!}{\begin{tabular}{lrrrrrrrr}
\toprule
时段 & \multicolumn{2}{c}{03-20} & \multicolumn{2}{c}{06-21} & \multicolumn{2}{c}{09-23} & \multicolumn{2}{c}{12-21} \\
\midrule
 & 原计划 & 调整后 & 原计划 & 调整后 & 原计划 & 调整后 & 原计划 & 调整后 \\
10:00--10:10 & 0.00 & 0.00 & 0.00 & 0.00 & 0.00 & 0.00 & 429.50 & 364.56 \\
12:00--12:10 & 571.96 & 501.26 & 0.00 & 0.00 & 624.53 & 0.00 & 965.03 & 951.01 \\
14:00--14:10 & 0.00 & 0.00 & 0.00 & 0.00 & 0.00 & 0.00 & 0.00 & 0.00 \\
16:00--16:10 & 692.18 & 512.10 & 108.82 & 108.82 & 613.48 & 609.17 & 791.37 & 823.78 \\
18:00--18:10 & 667.80 & 710.59 & 543.40 & 538.05 & 744.65 & 744.65 & 710.34 & 679.21 \\
20:00--20:10 & 0.00 & 0.00 & 0.00 & 0.00 & 0.00 & 0.00 & 0.00 & 0.00 \\
\bottomrule
\end{tabular}}
\end{table}
\begin{table}[htbp]\centering\scriptsize
\caption{指定日期各四小时区间的实际充放电量（kWh）}\label{tab:q3storage}
\resizebox{\textwidth}{!}{\begin{tabular}{lrrrrrrrr}
\toprule
时段 & \multicolumn{2}{c}{03-20} & \multicolumn{2}{c}{06-21} & \multicolumn{2}{c}{09-23} & \multicolumn{2}{c}{12-21} \\
\midrule
 & 充电 & 放电 & 充电 & 放电 & 充电 & 放电 & 充电 & 放电 \\
00:00--04:00 & 1478.05 & 183.06 & 116.38 & 54.80 & 2778.41 & 109.90 & 1716.84 & 79.29 \\
04:00--08:00 & 976.35 & 5858.46 & 58.68 & 3561.88 & 881.63 & 5997.91 & 2528.56 & 1896.21 \\
08:00--12:00 & 4366.52 & 2487.31 & 3767.71 & 0.00 & 5150.47 & 1896.62 & 5245.03 & 5659.53 \\
12:00--16:00 & 6171.66 & 254.72 & 6307.99 & 0.00 & 4602.61 & 520.92 & 5590.45 & 2220.11 \\
16:00--20:00 & 247.22 & 5398.14 & 10.35 & 4681.16 & 362.32 & 5546.48 & 11.08 & 5566.55 \\
20:00--24:00 & 3769.46 & 3034.39 & 8694.40 & 2485.90 & 7769.05 & 2974.46 & 7661.63 & 2937.45 \\
\bottomrule
\end{tabular}}
\end{table}
\begin{table}[htbp]\centering\small
\caption{指定日期连续紧急购电区间及电量}\label{tab:q3emergency}
\begin{tabular}{llr}
\toprule
日期 & 紧急购电区间 & 电量/kWh \\
\midrule
03-20 & 01:20--01:30 & 12.56 \\
 & 09:10--10:00 & 290.07 \\
 & 15:40--16:40 & 169.83 \\
 & 16:50--17:00 & 4.49 \\
06-21 & 无 & 0.00 \\
09-23 & 17:10--17:20 & 24.65 \\
 & 18:20--18:40 & 38.40 \\
 & 20:30--20:40 & 28.47 \\
 & 20:50--21:40 & 364.42 \\
12-21 & 无 & 0.00 \\
\bottomrule
\end{tabular}
\end{table}

主方案期末储电量为\QthreeEnd\ kWh。采用问题二相同的补充储电成本估值后，主方案相对0点基准的费用减少量为\QthreeInventoryDifference\ 元。完全信息费用下界仍为\QthreeLower\ 万元：对任意交付量，调整结算费用不低于按普通电价直接购买该交付量的费用，因此原完全信息问题仍提供有效下界。

为检验结算解释的影响，另将取消部分改为原计划照付并追加0.5倍违约费，完整重算全部预报方案，得到\QthreeLiteralCost\ 万元。在允许免费舍弃余电时，这一口径下减购通常没有优势：维持供给并舍弃多余电量即可避免新增违约费。因此，结算机制本身会改变调整方向，不能把不同机制下的结果直接归结为预报质量。

在四个指定日期固定历史信息与主方案日初状态，分别增加至14个情景，或改变远端储电目标，结果见表\ref{tab:q3sensitivity}。这些是同起点的单日检验，不用于重新选择全年主参数。
\begin{table}[htbp]\centering\scriptsize
\caption{问题三的单日敏感性比较（变化量相对主方案）}\label{tab:q3sensitivity}
\resizebox{\textwidth}{!}{\begin{tabular}{llrrrr}
\toprule
日期 & 设置 & 费用变化/元 & 紧急量变化/kWh & 日末储电/kWh & 库存校正变化/元 \\
\midrule
03-20 & 14情景 & -1565.19 & -246.57 & 4409.27 & -1467.16 \\
03-20 & 末端目标4800 & 0.00 & 0.00 & 4646.88 & 0.00 \\
03-20 & 末端目标7200 & 0.00 & 0.00 & 4646.88 & 0.00 \\
06-21 & 14情景 & 120.32 & 0.00 & 10800.00 & 67.05 \\
06-21 & 末端目标4800 & 0.00 & 0.00 & 10670.88 & 0.00 \\
06-21 & 末端目标7200 & 0.00 & 0.00 & 10670.88 & 0.00 \\
09-23 & 14情景 & -1297.26 & -189.31 & 8205.83 & -1313.59 \\
09-23 & 末端目标4800 & 0.00 & 0.00 & 8166.23 & 0.00 \\
09-23 & 末端目标7200 & 0.00 & 0.00 & 8166.23 & 0.00 \\
12-21 & 14情景 & 524.73 & 0.00 & 8508.92 & 420.61 \\
12-21 & 末端目标4800 & 0.00 & 0.00 & 8256.55 & 0.00 \\
12-21 & 末端目标7200 & 0.00 & 0.00 & 8256.55 & 0.00 \\
\bottomrule
\end{tabular}}
\end{table}

\subsection{检验与局限}
通过逐段复算能量平衡、储电量递推及交易分项，检查最终交付量与各时刻计划版本的一致性。改变未来实测或尚未发布的预报，不改变此前计划和执行动作；整点至十分钟的转换还通过常数功率、线性功率及跨午夜算例检验。全部策略的物理和费用残差均低于$10^{-5}$的相应单位容差。

十组全年实验均未触发维持计划或可行规则的求解失败回退。主方案出现一次求解状态异常，对同一线性规划冷启动重求解后获得最优解，实际执行仍采用优化结果。

新预报带来调整机会，但不保证每次实际费用都下降。有限历史情景可能遗漏突发偏差，线性插值忽略小时内更快的变化，远端补充预测和情景内储能前瞻也具有近似性。保持当前计划的候选检查保证的是同一模型中的预计费用关系；预报发布后的真实运行仍应由实际账单评价。本文仅比较附件提供的发布时间，不对未提供预报的其他时刻推断收益。
\begin{figure}[htbp]\centering
\includegraphics[width=\textwidth,height=.83\textheight,keepaspectratio]{figures/q3_day_20250320.pdf}
\caption{03-20的预报更新、购电调整与储能响应。}
\end{figure}

\clearpage
